\section{Conclusion}

We started from a simple game: $n$ people in a circle, every second person
eliminated, and the question of who is left at the end.

By playing the game by hand for $n=1,\dots,16$, we saw that person $1$ survives
when $n$ is a power of $2$, and that the survivor moves two places each time a
person is added. The key remark was that the game only depends on the places
of the people (Lemma~\ref{lem:places}), so that after one turn of the circle
the game is a smaller copy of itself. It gave the two relations
\[
J(2k)=2J(k)-1,\qquad J(2k+1)=2J(k)+1,
\]
and, by induction on $m$, the formula
\[
J(2^m+\ell)=2\ell+1,\qquad 0\leq\ell<2^m ,
\]
which we proved a second time by induction on $n$. Since $2^m$ is the largest
power of $2$ not bigger than $n$, it can also be written
$J(n)=2n-2^{m+1}+1$ with $m=\lfloor\log_2 n\rfloor$.
We then removed the assumption that person $1$ starts: if the game starts with
person $s$, the survivor is the person $2\ell$ places after $s$. In base $2$,
the formula becomes a rule on the digits: the first digit of $n$ moves to the
end. This gives the survivor of a circle of $10^9$ people at once, where the
game would need almost a billion eliminations. Every result was checked by a
computer simulation.

The work followed the method of the course: test small cases, look for a
pattern, state it precisely, and prove it. Each step was needed. The tests
gave the formula, but they could not prove it: sixteen cases say nothing about
the seventeenth. The proof, in turn, came from looking closely at one round of
the game in the tests. The computer was useful at both ends, first to check
our drawings, then to check the formulas on many more cases than we could do
by hand. It also showed us that a formula is only fast if it is used well:
our first program based on the formula was slower than the simulation.

Some questions remain open for us. The first one is to eliminate every third
person instead of every second one, or more generally every $q$-th person. The
proof of Lemma~\ref{lem:places} still works: one step now sends the first
$q-1$ people to the back of the list and eliminates the next one. It gives
\[
J_q(n)=\big((J_q(n-1)+q-1)\bmod n\big)+1,
\]
where $a\bmod n$ is the remainder of $a$ divided by $n$. For $q=2$ this is
relation~\eqref{eq:onestep}. We checked this relation with a simulation for
$q=3,4,5$ and $n<150$. For $q=3$ the survivors for $n=1,\dots,16$ are
$1,2,2,1,4,1,4,7,1,4,7,10,13,2,5,8$: the survivor goes up by $3$ and comes
back near the start from time to time, but we did not find a formula as
simple as $2\ell+1$. The case $q=2$ is special because the circle is exactly
halved at each round. A last question remains: can we predict the whole
order of the eliminations, and not only the survivor?
